\section{An explicit constructive separation of polynomial and exponential time}
\label{sec:diagonal}

The foundation should support lower bounds as well as the verification of
upper bounds.  We now give a direct separation
\[
  \mathrm{P}\subsetneq\mathrm{EXPTIME}
\]
whose diagonal counterexample is a primitive-recursive function of the
proposed polynomial algorithm.  The argument deliberately uses a generous
exponential upper bound.  It does not require the sharp time hierarchy theorem,
an asymptotic domination principle, or any claim about
\(\mathrm{P}\) versus \(\mathrm{NP}\).

\subsection{The exact bounded simulation convention}

Here the source model is a deterministic one-tape machine with a fixed finite
tape alphabet containing the two input symbols and a blank.  The input word
occupies consecutive cells beginning at position zero; the head starts at
zero.  Accepting and rejecting states are absorbing.  Every binary word is
allowed as a program description: a word that does not describe a legal
transition table denotes a fixed machine that rejects immediately.  In
particular the bounded acceptance function is total for every description.
Here \(e\) ranges over binary-word codes \(e\ge1\), including malformed
machine payloads. A numerical descriptor containing 0 is normalized to
the empty description before this construction, as in the machine convention.

Write
\[
  B(e,w,t)\in\{0,1\}
\]
for acceptance of program \(e\) on input \(w\) after \(t\) iterations of its
transition function.  Absorption makes this equivalent to acceptance by time
\(t\).  It is a primitive-recursive function of all three natural-number
codes.  No assertion that the original program eventually halts is involved.
The code of a polynomial presentation is a triple \((e,c,k)\), where
\(c\geq 1\) and \(k\geq0\), with Boolean interpretation
\begin{equation}
  L_{e,c,k}(w)=B\bigl(e,w,c(|w|+1)^k\bigr).
  \label{eq:diag-presentation}
\end{equation}
Every such presentation defines a polynomial-time language, since a bounded
universal simulation incurs only a polynomial overhead.  Conversely, every
ordinary polynomial-time decision algorithm has a presentation of this form,
after enlarging its constant and degree.  A presentation that times out still
has a defined answer: its value in \eqref{eq:diag-presentation} is zero.

\begin{lemma}[A coarse physical simulation bound]
\label{lem:coarse-universal}
There is a fixed one-tape simulator and a fixed positive integer \(C\) that
compute \(B(e,w,t)\), supplied with a unary countdown \(t\), in at most
\[
  C\bigl(|e|+|w|+t+1\bigr)^4
\]
actual tape transitions.  The simulator, its correctness proof, and this
bound admit primitive-recursive descriptions and quantifier-free induction
proofs in iPRA.
\end{lemma}

\begin{proof}
First describe a simulator on a fixed finite number of virtual tapes.  One
tape holds the simulated tape; another holds the source transition table;
others hold the binary source-state code, scan markers, and the countdown.
If \(m=|e|\), a deliberately inefficient full scan of the transition table,
using repeated scans to compare state codes, costs at most a fixed multiple
of \((m+1)^2\) virtual transitions per simulated move.  The source-state code
has at most a fixed multiple of \(m+1\) bits under a conventional transition
table encoding.  A malformed table is detected by a finite parsing scan and
replaced by the rejecting program.

During \(t\) simulated moves, the simulated head can visit only the interval
\([-t,|w|+t]\).  The virtual tapes therefore occupy at most a fixed multiple
of
\(N=m+|w|+t+1\) cells, including the unary countdown.  Encode their cells on
one physical tape by a fixed number of tracks, with a marked virtual head on
each track.  A bounded number of complete sweeps reads the marked symbols,
updates them, and moves their markers.  Each virtual move costs at most a
fixed multiple of \(N\) physical moves.  Thus the main simulation costs at
most a fixed multiple of
\[
  t(m+1)^2N\leq N^4.
\]
Initialization, parsing, copying, and clearing markers fit the same bound
after increasing the fixed constant.

For formal verification, the invariant says that every encoded virtual cell
and every head marker agrees with the corresponding simulated configuration,
that the countdown equals the remaining fuel, and that the occupied interval
has the displayed bound.  All cell comparisons and scan-prefix conditions
are bounded by explicit numerical terms.  Their Boolean conjunction is a
quantifier-free induction invariant.  The running-time argument uses the
same invariant together with a transition counter.  The constant \(C\) is
fixed by the chosen finite scan implementation; this proof does not pretend
to supply an optimized numerical transition count.
\end{proof}

\subsection{A description that carries its own clock}

Let \(1^r\) denote a unary block of \(r\) ones, and let juxtaposition denote
word concatenation.  Define
\begin{equation}
  W(e,c,k)=1^{|e|}\,0\,e\,1^c\,0\,1^k\,0,
  \qquad e\ge1,\ c\geq1.
  \label{eq:self-description}
\end{equation}
The first zero terminates the unary length of \(e\).  The parser then reads
exactly that many bits as the program payload, so zeros inside \(e\) cause no
ambiguity.  The remaining two unary fields end at the next two zeros; the
word must end there.  The field for \(k=0\) consists just of its terminating
zero.  A payload need not be a legal machine description, because the default
rejecting interpretation has already made the semantics total.

The constructor \(W\) and this parser are primitive recursive.  Direct
concatenation identities prove that parsing \(W(e,c,k)\) returns exactly
\((e,c,k)\).  Its length is
\begin{equation}
  n=|W(e,c,k)|=2|e|+c+k+3.
  \label{eq:diag-length}
\end{equation}
Consequently \(|e|\leq n\), \(c\leq n\), and \(k\leq n\).

Define the Boolean language \(D\) by the following finite algorithm.  On
input \(w\), parse the displayed grammar.  If parsing fails, or if the middle
unary field has length zero, return zero.  Otherwise recover \(e,c,k\), put
\(n=|w|\), and return
\begin{equation}
  D(w)=1-B\bigl(e,w,c(n+1)^k\bigr).
  \label{eq:diag-definition}
\end{equation}
The input \(w\) itself is passed to the simulated program.  All operations
in \eqref{eq:diag-definition} are primitive recursive: parsing, length, exponentiation, bounded
iteration, and Boolean complementation.  This establishes definability in
iPRA.  A separate machine-cost proof is needed to establish the asserted
exponential upper bound.

\begin{lemma}[An exponential bound for the diagonal language]
\label{lem:diagonal-exp}
For a fixed positive integer \(K\), a fixed deterministic machine computes
\(D(w)\) in at most
\[
  2^{K(|w|+1)^2}
\]
transitions on every input \(w\).  The bound and correctness have iPRA proofs.
\end{lemma}

\begin{proof}
Consider a well-formed input of length \(n\).  By \eqref{eq:diag-length}, its recovered constants
satisfy \(1\leq c\leq n\), \(0\leq k\leq n\), and \(n\geq1\).  Set
\(T=c(n+1)^k\).  Quantifier-free induction proves \(n+1\leq2^n\) for
\(n\geq1\), whence
\begin{equation}
  T\leq n(n+1)^n
    \leq 2^n 2^{n^2}
    =2^{n^2+n}
    \leq 2^{(n+1)^2}.
  \label{eq:diag-clock-bound}
\end{equation}

An especially simple implementation produces a unary countdown of length
\(T\).  Start with the unary \(c\) block and multiply its length by \(n+1\),
\(k\) times, by repeated copying.  Every intermediate length is at most
\(T\).  In a naive tape implementation, a transferred unary mark requires
only a fixed number of scans of the input, old block, new block, and loop
markers, occupying \(O(n+T+1)\) cells.  There are at most \(n\) multiplication
stages and at most \(T\) transferred marks per stage, with polynomially
bounded loop overhead.  Thus the generous bound
\(A_0(n+T+1)^4\) covers countdown production.  The simulation lemma, with
\(|e|\leq n\), covers evaluation of \(B\) by another fixed multiple of the
same fourth power.  Parsing malformed inputs and returning zero costs only
a polynomial, also covered after increasing the constant.  For some fixed
positive integer \(A\), the complete algorithm therefore has cost
\[
  A(n+T+1)^4
\]
on well-formed inputs, with the polynomial parsing bound on the others.

Put \(q=(n+1)^2\).  Both \(T\) and \(n+1\) are at most \(2^q\), so
\[
  A(n+T+1)^4\leq 16A\,2^{4q}\leq2^{Kq}
\]
for any fixed \(K\geq4+\lceil\log_2(16A)\rceil\).  If one wishes to avoid
logarithms even in the specification of the constant, \(K=4+16A\) works.
Increase \(A\), and hence \(K\), to cover invalid inputs and the finitely many
small input lengths.  These are explicit polynomial and exponential bounds,
not an appeal to an eventual asymptotic comparison.

The iPRA audit uses Boolean prefix invariants for parsing and copying, the
simulator invariant of Lemma~\ref{lem:coarse-universal}, and the elementary
quantifier-free inductions for multiplication, exponentiation, and \eqref{eq:diag-clock-bound}.
All storage and all loop lengths are bounded by the displayed
primitive-recursive terms.  Unary countdown production is intentionally
inefficient, but its cost remains within the claimed exponential bound.
\end{proof}

\begin{theorem}[A uniform primitive-recursive mismatch witness]
\label{thm:explicit-diagonal}
iPRA proves that, for every program description \(e\), every \(c\geq1\), and
every \(k\geq0\), with description code \(e\ge1\),
\begin{equation}
  D\bigl(W(e,c,k)\bigr)
  \ne
  B\!\left(e,W(e,c,k),
       c\bigl(|W(e,c,k)|+1\bigr)^k\right).
  \label{eq:diag-mismatch}
\end{equation}
Consequently the displayed \(D\) belongs to \(\mathrm{EXPTIME}\) and has no
polynomial-time presentation.  In the intended extensional interpretation
of the coded classes, \(\mathrm{P}\subsetneq\mathrm{EXPTIME}\).
\end{theorem}

\begin{proof}
The parser equations give the exact triple \((e,c,k)\) on input
\(W(e,c,k)\).  Substituting it in \eqref{eq:diag-definition}, the left side of \eqref{eq:diag-mismatch} is the Boolean
complement of the right side.  The two possible Boolean values show that
they are unequal.  This uses no totality promise for the original program:
both rejecting and nonhalting computations have bounded-acceptance value
zero, whose complement is one.

Lemma~\ref{lem:diagonal-exp} supplies the exponential-time upper bound.
Equation \eqref{eq:diag-mismatch} rules out each polynomial presentation, with the actual
counterexample word given by the same primitive-recursive constructor
\(W\).  If an ordinary polynomial-time decider computed \(D\), choose a
clock \(c(n+1)^k\) dominating its running time.  Its bounded-acceptance
value would equal \(D\) on every word, contradicting \eqref{eq:diag-mismatch} on
\(W(e,c,k)\). For the non-strict inclusion, replace each polynomial
descriptor by the total clocked evaluator that returns its Boolean answer,
including its timeout answer. That evaluator has a genuine polynomial
runtime bound by the simulator lemma. An exponential clock dominates this
bound and therefore preserves its answer. Merely increasing the clock of
the original possibly partial program would not justify the inclusion.
\end{proof}

\subsection{What the separation does, and what its formalization must retain}

Equation \eqref{eq:diag-mismatch} is stronger, constructively, than merely negating the existence
of a polynomial decider.  Intuitionistic logic does not in general turn
\(\neg\forall w\,R(w)\) into \(\exists w\,\neg R(w)\).  Here there is no such
conversion: \(W(e,c,k)\) explicitly gives a differing input for each proposed
presentation.  The statement can be expressed as the universal closure of
a quantifier-free Boolean inequality with named primitive-recursive
functions.  The required constructor/parser identities are themselves
proved by quantifier-free induction on finite word operations.

The unary clock fields are essential to the resource estimate.  If the
exponent \(k\) were given in unrestricted binary and then used without a
length-dependent cap, it could be exponential in \(|w|\), and the same
algorithm could require more than exponential time.  Unary encoding still
represents every fixed polynomial clock, while guaranteeing \(k,c\leq|w|\)
on the specially constructed inputs and every successfully parsed input.

The argument also keeps program evaluation and logical reflection separate.
It evaluates a specified machine for a specified finite number of steps;
it does not evaluate arbitrary primitive-recursive definitions uniformly,
decide whether an arbitrary theorem is true, or infer soundness from an
arbitrary iPRA proof code.  The familiar diagonal phenomenon therefore poses
no conflict with the totality of the bounded evaluator.

For a Lean formalization, the key interfaces are the word constructor,
its exact parsing theorem, the bounded evaluator, its Boolean-valuedness,
the simulator's physical transition bound, and the elementary exponential
inequalities.  The final mismatch proof is short once those interfaces have
been established.  A checked Lean proof of the mismatch algebra alone does
not establish the machine-time bound or automatically produce an iPRA
derivation.  The completed development must retain the finite invariants and
the code-generating bridge described above for those separate obligations.

